\section{Operational Appendix: End-to-End Execution Recipe}

This appendix provides an explicit command-level runbook so the complete artifact can be regenerated without ambiguity. The intent is to make grading and independent verification straightforward, even if the notebook is executed on a fresh machine with only the repository and Python environment available.

\subsection{Command-Level Runbook}

\begin{lstlisting}[basicstyle=\ttfamily\scriptsize,frame=single]
# 1) Activate environment and move to project root
source /Users/tahamajs/Documents/uni/venv/bin/activate
cd /Users/tahamajs/Documents/uni/NLP/nlp-assignments-spring-2023/NLP_UT/CA2

# 2) Verify key files
ls answer/NLP_HW2.ipynb
ls datasets/q1/emails_1.csv
ls datasets/q1/emails_2.csv
ls datasets/q2/urls.csv

# 3) Ensure required Python dependencies are present
python -c "import pandas, numpy, sklearn, matplotlib, seaborn, tldextract"

# 4) Execute notebook non-interactively
MPLBACKEND=Agg jupyter nbconvert --to notebook --execute answer/NLP_HW2.ipynb \
  --output /tmp/ca2_exec_check.ipynb --ExecutePreprocessor.timeout=3600

# 5) Promote executed notebook as final deliverable
cp /tmp/ca2_exec_check.ipynb answer/NLP_HW2.ipynb

# 6) Validate expected figure outputs
ls answer/report/figures/q1_confusion_matrices.png
ls answer/report/figures/q1_roc_pr_curves_lr.png
ls answer/report/figures/q1_top_words_spam_ham.png
ls answer/report/figures/q1_model_comparison_bar.png
ls answer/report/figures/q2_structural_distributions.png
ls answer/report/figures/q2_correlation_heatmap_all_features.png
ls answer/report/figures/q2_structural_boxplots.png
ls answer/report/figures/q2_pairplot_sample.png
ls answer/report/figures/q2_top_mean_abs_correlations.png
ls answer/report/figures/q2_roc_combined_features.png
ls answer/report/figures/q2_pr_combined_lr.png
ls answer/report/figures/q2_lr_feature_importance.png

# 7) Build IEEE report
cd answer/report
pdflatex -interaction=nonstopmode main.tex
bibtex main
pdflatex -interaction=nonstopmode main.tex
pdflatex -interaction=nonstopmode main.tex

# 8) Check final page count and PDF metadata
pdfinfo main.pdf | rg '^Pages|^Title'

# 9) Optional quality checks
# - Open main.pdf and visually inspect all figure captions and references
# - Confirm no missing figures in log
# - Confirm model tables match notebook outputs

# 10) Final artifact set
# - answer/NLP_HW2.ipynb
# - answer/report/main.pdf
# - answer/report/figures/*.png
# - answer/report/main.tex + section files + references.bib
\end{lstlisting}

\subsection{Expected Observable Outcomes}

A successful run should satisfy the following observable outcomes:
\begin{enumerate}
    \item Notebook execution terminates without \texttt{CellExecutionError}.
    \item Exactly twelve report figures are present and non-empty.
    \item The report compiles with IEEEtran and includes all assignment sections.
    \item Final PDF page count is strictly greater than ten pages.
\end{enumerate}

This appendix intentionally duplicates procedural information from earlier sections in a command-first format, because reproducibility failures in assignments frequently occur in operational details rather than in modeling logic.
